whereas a full categorical treatment would insist on
arbitrary tensor products. In fact, more general notions
already exist in the literature. Our planar algebras could
be formulated as a rather special kind of ``spider" in
the sense of Kuperberg in [Ku], or one could place them in
the context of pivotal and spherical categories ([FY],[BW]), 
and the theory of $C^*-$tensor categories even has the 
ever desirable positivity ([LR],[We3]). Also, in the semisimple
case at least, the work in section 3.3 on cabling and reduction 
shows how to extend our planar diagrams to ones with labelled
edges.

But it is the very restrictive nature of our definition of planar
algebras that should be its great virtue. We have good reasons
for limiting the generality. The most compelling is the 
equivalence with subfactors, which has been our guiding light.
We have tried to introduce as little formalism as possible
compatible with exhibiting quite clearly the planar nature of
subfactor theory. Thus our intention has been to give pride
of place to the pictures. But subfactors are not the only
reason for our procedure. By restricting the scope of the theory
one hopes to get to the most vital examples as quickly as
possible. And we believe that we will see, in some form,
{\it all} the examples in our restricted theory anyway. 
Thus the Fuss-Catalan algebras of [BJ2] (surely among the 
most basic planar algebras, whatever one's definition)
first appeared with our strict axioms. Yet at the same time,
as we show in section 2.5, the {\sc homfly} polynomial, for
which one might have thought oriented strings essential, can
be completely captured within our unoriented framework.

It is unlikely that any other restriction of some more general
operad is as rich as the one we use here. To see why, note that 
in the operadic picture, the role of the identity is played by 
tangles without internal discs -see [Ma]. In our case we get the
whole Temperley-Lieb algebra corresponding to the identity 
whereas any orientation restriction will reduce the size of
this ``identity". The beautiful structure of the Temperley-Lieb
algebra is thus always at our disposal. This leads to the 
following rather telling reason for looking carefully 
at our special planar algebras among more general ones:
if we introduce the generating function for the dimensions
of a planar algebra,  $\sum^\i_{n=0}\dim (V_n)z^n$,
we shall see that if the planar algebra satisfies reflection 
positivity, then this power series has non-zero radius of
convergence. By contrast, if we take the natural oriented planar
algebra structure given by the {\sc homfly} skein, it is a 
result of Ocneanu and Wenzl ([We1],[F+]) that there is a positive
definite Markov trace on the whole algebra, even though the
generating function has zero radius of convergence.

In spite of the previous polemic, it would be foolish to 
neglect the fact that our planar algebra formalism fits
into a more general one. Subfactors can be constructed
with arbitrary orientations by the procedure of [We1],[F+]
and it should be possible to calculate their planar algebras
by planar means.

We end this introduction by discussing three of our motivations
for the introduction of planar algebras as we have defined them.

{\it Motivation 1 }  Kauffman gave his now well-known pictures
for the Temperley-Lieb algebra in [Ka1]. In the mid 1980's he 
asked the author if it was possible to give a pictorial
representation of all elements in the tower of algebras of [J1].
We have only developed the planar algebra formalism for the
sub-tower of relative commutants, as the all-important rotation is
not defined on the whole tower. Otherwise this paper constitutes
an answer to Kauffman's question.

{\it Motivation 2 } One of the most extraordinary developments in 
subfactors was the discovery by Haagerup in [Ha] of a subfactor
of index $(5 + \sqrt{13})/2$, along with the proof that this is 
the smallest index value, greater than $4$, of a finite depth
subfactor. As far as we know there is no way to obtain Haagerup's
``sporadic" subfactor from the conformal field theory/quantum group
methods of [Wa],[We3],[X],[EK]. It is our hope that the planar algebra
context will put Haagerup's subfactor in at least one natural
family, besides yielding tools for its study that are more
general than those of [Ha]. For instance it follows from 
Haagerup's results that the planar algebra of his subfactor is
generated by a single element in $V_4$ (a ``4-box"). The small
dimensionality of the planar algebra forces extremely strong 
conditions on this 4-box. The only two simpler such planar algebras
(with reflection positivity) are those of the $D_6$ subfactor
of index $4\cos^2\pi/10$ and the $\tilde E_7$ subfactor of
index $4$. There are analogous planar algebras generated by
2-boxes and 3-boxes. The simplest 2-box case comes from the
$D_4$ subfactor (index $3$) and the two simplest 3-box cases from
$E_6$ and $\tilde E_6$ (indices $4\cos^2\pi/12$ and $4$).
Thus we believe there are a handful of planar algebras for
each k, generated by a single k-box, satisfying extremely 
strong relations. Common features among these relations
should yield a unified calculus for constructing and manipulating
these planar algebras. In this direction we have classified 
with Bisch in [BJ1] {\it all} planar algebras generated by
a 2-box and tightly restricted in dimension. A result of
D.Thurston shows an analogous result should exist for 3-boxes - see section 2.5.
The 4-box case has yet to be attempted.

In general one would like to understand all systems of relations
on planar algebras that cause the free planar algebra to collapse
to finite dimensions. This is out of sight at the moment. 
Indeed it is know from [BH] that subfactors of index $6$ are 
``wild" in some technical sense, but up to $3+\sqrt{3}$ they
appear to be ``tame". It would be significant to know for what   
index value subfactors first become wild.

{\it Motivation 3 } Since the earliest days of subfactors it
has been known that they can be constructed from certain 
finite data known as a commuting square (see [GHJ]). A theorem 
of Ocneanu (see [JS] or [EK]) reduced the problem of  
calculating the planar algebra component $V_k$ of such a subfactor
to the solution of a finite system of linear equations in finitely
many unknowns. Unfortunately the number of equations 
grows exponentially with $k$ and it is unknown at present whether 
the most simple questions concerning these $V_k$ are solvable
in polynomial time or not. On the other hand the planar algebra
gives interesting invariants of the original combinatorial 
data and it was a desire to exploit this information that led
us to consider planar algebras. First it was noticed that there
is a suggestive planar notation for the linear equations themselves.
Then the invariance of the solution space under the action
of planar tangles was observed. It then became clear that 
one should consider other ways of constructing planar 
algebras from combinatorial data, such as the planar
algebra generated by a tensor in the tensor planar algebra.

These ideas were the original motivation for introducing planar
algebras. We discuss these matters in more detail in section
2.11 which is no doubt the most important part of this work.
The significance of Popa's result on $\lambda -$ lattices became
apparent as the definition evolved. Unfortunately we have
not yet been able to use planar algebras in a convincing
way as a tool in the calculation of the planar algebra for
specific commuting squares.

This paper has been written over a period of several years and
many people have contributed. In particular I would like to
thank Dietmar Bisch, Pierre de la Harpe, Roland Bacher, Sorin Popa,
Dylan Thurston, Bina Bhattacharya, Zeph Landau, Adrian Ocneanu, Gib Bogle and 
Richard Borcherds. Deborah Craig for her 
patience and first-rate typing, and Tsukasa Yashiro for the pictures. 


  
\bs\bs
\ni{\large\bf 1. The Formalism}

\medskip
{\bf Definition 1.1.} If $k$ is a non-negative integer, the
{\it standard} $k${\it -box}, $\b _k$, is
$\{ (x,y)\in\r\mid 0\leq x\leq k\!+\!1, \
0\leq y\leq 1\}$, together with the $2k$ marked points,
$1\!=\!(1,1)$, $2\!=\!(2,1)$, $3=(3,1),\dots ,k=(k,1)$,
$k+1=(k,0)$, \ $k+2=(k-1,0),\dots ,2k=(1,0)$.

\bs{\bf Definition 1.2.}
A {\it planar network} $\cal N$ will be a subset of $\r$ consisting
of the union of a finite set of disjoint images of $\b_k$'s
(with $k$ varying) under smooth orientation-preserving
diffeomorphisms of $\r$, and a finite number of oriented disjoint
curves, smoothly embedded, which may be closed (i.e. isotopic to
circles), but if not their endpoints coincide with marked points of
the boxes. Otherwise the curves are disjoint from the boxes.
All the marked points are endpoints of curves, which meet the boxes
transversally. The orientations of the curves must satisfy the
following two conditions.
\begin{itemize}
\item[a)] A curve meeting a box at an odd marked point
must exit the box at that point.
\item[b)] The connected components of $\r\backslash\cal N$
may be oriented in such a way that the orientation of a curve
coincides with the orientation induced as part of the boundary of a
connected component.
\end{itemize}

{\bf Remark.} Planar networks are of two kinds according to the
orientation of the unbounded region.

\bs
Let $L_i$, $i=1,2,\dots$ be sets and $L=\coprod_i L_i$ be their
disjoint union. $L$ will be called the set of ``labels".

\bs{\bf Definition 1.3.} A {\it labelled planar network}
(on $L$) will be a planar network together with a function from its
$k$-boxes to $L_k$, for all $k$ with $L_k\neq\emptyset$.

\bs If the labelling set consists of asymmetric letters, we may
represent the labelling function diagrammatically by placing the
corresponding letter in its box, with the understanding that the
first marked point is at the top left. This allows us to ignore the
orientations on the edges and the specification of the marked
points. In Fig.~1.4 we give an example of a labelled planar network
with $L_1=\{ P\}$, $L_2=\{ R\}$, $L_3=\{ Q\}$. Here the unbounded
region is positively oriented and, in order to make the conventions
quite clear, we have explicitly oriented the edges and numbered the
marked points of the one 3-boxes labelled $Q$.
\[
        \begin{picture}(0,0)%
\epsfig{file=ps/pic7.ps}%
\end{picture}%
\setlength{\unitlength}{0.00041700in}%
%
\begingroup\makeatletter\ifx\SetFigFont\undefined
% extract first six characters in \fmtname
\def\x#1#2#3#4#5#6#7\relax{\def\x{#1#2#3#4#5#6}}%
\expandafter\x\fmtname xxxxxx\relax \def\y{splain}%
\ifx\x\y   % LaTeX or SliTeX?
\gdef\SetFigFont#1#2#3{%
  \ifnum #1<17\tiny\else \ifnum #1<20\small\else
  \ifnum #1<24\normalsize\else \ifnum #1<29\large\else
  \ifnum #1<34\Large\else \ifnum #1<41\LARGE\else
     \huge\fi\fi\fi\fi\fi\fi
  \csname #3\endcsname}%
\else
\gdef\SetFigFont#1#2#3{\begingroup
  \count@#1\relax \ifnum 25<\count@\count@25\fi
  \def\x{\endgroup\@setsize\SetFigFont{#2pt}}%
  \expandafter\x
    \csname \romannumeral\the\count@ pt\expandafter\endcsname
    \csname @\romannumeral\the\count@ pt\endcsname
  \csname #3\endcsname}%
\fi
\fi\endgroup
\begin{picture}(5676,5943)(1164,-6531)
\put(1366,-2951){\makebox(0,0)[lb]{\smash{\SetFigFont{12}{14.4}{rm}$P$}}}
\put(2836,-4106){\makebox(0,0)[lb]{\smash{\SetFigFont{12}{14.4}{rm}$R$}}}
\put(4891,-5591){\makebox(0,0)[lb]{\smash{\SetFigFont{12}{14.4}{rm}$R$}}}
\put(5618,-2689){\makebox(0,0)[lb]{\smash{\special{ps:gsave currentpoint currentpoint translate
-135.0 rotate neg exch neg exch translate}\SetFigFont{12}{14.4}{rm}$Q$\special{ps:currentpoint grestore moveto}}}}
\put(6451,-2911){\makebox(0,0)[lb]{\smash{\special{ps:gsave currentpoint currentpoint translate
-135.0 rotate neg exch neg exch translate}\SetFigFont{12}{14.4}{rm}6\special{ps:currentpoint grestore moveto}}}}
\put(5476,-1936){\makebox(0,0)[lb]{\smash{\special{ps:gsave currentpoint currentpoint translate
-135.0 rotate neg exch neg exch translate}\SetFigFont{12}{14.4}{rm}4\special{ps:currentpoint grestore moveto}}}}
\put(5851,-2311){\makebox(0,0)[lb]{\smash{\special{ps:gsave currentpoint currentpoint translate
-135.0 rotate neg exch neg exch translate}\SetFigFont{12}{14.4}{rm}5\special{ps:currentpoint grestore moveto}}}}
\put(5776,-3436){\makebox(0,0)[lb]{\smash{\special{ps:gsave currentpoint currentpoint translate
-135.0 rotate neg exch neg exch translate}\SetFigFont{12}{14.4}{rm}1\special{ps:currentpoint grestore moveto}}}}
\put(5326,-2986){\makebox(0,0)[lb]{\smash{\special{ps:gsave currentpoint currentpoint translate
-135.0 rotate neg exch neg exch translate}\SetFigFont{12}{14.4}{rm}2\special{ps:currentpoint grestore moveto}}}}
\put(4876,-2461){\makebox(0,0)[lb]{\smash{\special{ps:gsave currentpoint currentpoint translate
-135.0 rotate neg exch neg exch translate}\SetFigFont{12}{14.4}{rm}3\special{ps:currentpoint grestore moveto}}}}
\end{picture}

\]
\begin{center}
Figure 1.4
\end{center}

\bs\ni
Note that the same picture as in Fig.~1.4, but with an $R$ upside
down, would be a different labelled planar tangle since the marked
points would be different.
With or without labels, it is only  necessary to say which
distinguished boundary point is first. 

\bs{\bf Remark 1.5.} By shrinking each $k$-box to a point as in
Fig.~1.6 one obtains from a planar network a system of immersed
curves with transversal multiple point singularities.
\[
        \begin{picture}(0,0)%
\epsfig{file=ps/pic8.ps}%
\end{picture}%
\setlength{\unitlength}{0.00041700in}%
%
\begingroup\makeatletter\ifx\SetFigFont\undefined
% extract first six characters in \fmtname
\def\x#1#2#3#4#5#6#7\relax{\def\x{#1#2#3#4#5#6}}%
\expandafter\x\fmtname xxxxxx\relax \def\y{splain}%
\ifx\x\y   % LaTeX or SliTeX?
\gdef\SetFigFont#1#2#3{%
  \ifnum #1<17\tiny\else \ifnum #1<20\small\else
  \ifnum #1<24\normalsize\else \ifnum #1<29\large\else
  \ifnum #1<34\Large\else \ifnum #1<41\LARGE\else
     \huge\fi\fi\fi\fi\fi\fi
  \csname #3\endcsname}%
\else
\gdef\SetFigFont#1#2#3{\begingroup
  \count@#1\relax \ifnum 25<\count@\count@25\fi
  \def\x{\endgroup\@setsize\SetFigFont{#2pt}}%
  \expandafter\x
    \csname \romannumeral\the\count@ pt\expandafter\endcsname
    \csname @\romannumeral\the\count@ pt\endcsname
  \csname #3\endcsname}%
\fi
\fi\endgroup
\begin{picture}(6674,2204)(1179,-4023)
\end{picture}

\]
\begin{center}
Figure 1.6
\end{center}
\ni
Cusps can also be handled by labelled 1-boxes. To reverse the
procedure requires a choice of incoming curve at each multiple
point but we see that our object is similar to that of Arnold in
[A]. In particular, in what follows we will construct a huge
supply of invariants for systems of immersed curves. It remains to
be seen whether these invariants are of interest in singularity
theory, and whether Arnold's invariants may be used to construct
planar algebras with the special properties we shall describe.

\bs{\bf Definition 1.7.} A planar $k$-tangle $T$ (for
$k=0,1,2,\dots$) is the intersection of a planar network $\cal N$
with the standard $k$-box $\b_k$, with the condition that the
boundary of $\b_k$ meets $\cal N$ transversally precisely in the
set of marked points of $\b_k$, which are points on the curves of
$\cal N$ other than endpoints. The orientation induced by $\cal N$
on a neighborhood of (0,0) is required to be positive. A labelled
planar $k$-tangle is defined in the obvious way.

The connected curves in a tangle $T$ will be called the {\it
strings} of $T$. 

The set of smooth isotopy classes of labelled planar $k$-tangles,
with isotopies being the identity on the boundary of $\b_k$, is
denoted $T_k(L)$.

{\bf Note.} $T_0(L)$ is naturally identified with the set of planar
isotopy classes of labelled networks with unbounded region
positively oriented.

\bs{\bf Definition 1.8.} The associative algebra
$\P_k(L)$ over the field $K$ is the vector space having $T_k(L)$ as
basis, with multiplication defined as follows. If $T_1,T_2\in
T_k(L)$, let
$\tilde T_2$ be $T_2$ translated in the negative $y$ direction by
one unit. After isotopy if necessary we may suppose that the union
of the curves in $T_1$ and $\tilde T_2$ define smooth curves.
Remove  $\{(x,0)\mid 0\leq x\leq k+1, \ x\not\in\z\}$ from
$T_1\cup\tilde T_2$ and finally rescale by multiplying the
$y$-coordinates $\frac 12$, then adding $\frac 12$.
The resulting isotopy class of labelled planar $k$-tangles is
$T_1T_2$. See Figure 1.9 for an example.
\[
        \begin{picture}(0,0)%
\epsfig{file=ps/pic9.ps}%
\end{picture}%
\setlength{\unitlength}{0.00041700in}%
%
\begingroup\makeatletter\ifx\SetFigFont\undefined
% extract first six characters in \fmtname
\def\x#1#2#3#4#5#6#7\relax{\def\x{#1#2#3#4#5#6}}%
\expandafter\x\fmtname xxxxxx\relax \def\y{splain}%
\ifx\x\y   % LaTeX or SliTeX?
\gdef\SetFigFont#1#2#3{%
  \ifnum #1<17\tiny\else \ifnum #1<20\small\else
  \ifnum #1<24\normalsize\else \ifnum #1<29\large\else
  \ifnum #1<34\Large\else \ifnum #1<41\LARGE\else
     \huge\fi\fi\fi\fi\fi\fi
  \csname #3\endcsname}%
\else
\gdef\SetFigFont#1#2#3{\begingroup
  \count@#1\relax \ifnum 25<\count@\count@25\fi
  \def\x{\endgroup\@setsize\SetFigFont{#2pt}}%
  \expandafter\x
    \csname \romannumeral\the\count@ pt\expandafter\endcsname
    \csname @\romannumeral\the\count@ pt\endcsname
  \csname #3\endcsname}%
\fi
\fi\endgroup
\begin{picture}(13006,3952)(511,-5178)
\put(511,-3356){\makebox(0,0)[lb]{\smash{\SetFigFont{12}{14.4}{rm}$T_1=$}}}
\put(4036,-3431){\makebox(0,0)[lb]{\smash{\SetFigFont{12}{14.4}{rm},}}}
\put(4426,-3401){\makebox(0,0)[lb]{\smash{\SetFigFont{12}{14.4}{rm}$T_2=$}}}
\put(7966,-3446){\makebox(0,0)[lb]{\smash{\SetFigFont{12}{14.4}{rm},}}}
\put(8341,-3476){\makebox(0,0)[lb]{\smash{\SetFigFont{12}{14.4}{rm}$T_1T_2=$}}}
\put(2531,-3266){\makebox(0,0)[lb]{\smash{\SetFigFont{12}{14.4}{rm}$Q$}}}
\put(11387,-2187){\makebox(0,0)[lb]{\smash{\SetFigFont{12}{14.4}{rm}$Q$}}}
\put(10826,-4527){\makebox(0,0)[lb]{\smash{\SetFigFont{12}{14.4}{rm}$R$}}}
\put(6176,-3401){\makebox(0,0)[lb]{\smash{\SetFigFont{12}{14.4}{rm}$R$}}}
\end{picture}

\]
\begin{center}
        Figure 1.9
\end{center}
\bs{\bf Remark.} The algebra $\P_k(L)$ has an obvious unit and
embeds unitally in $\P_{k+1}(L)$ by adding the line
$\{(k+1,t)\mid 0\leq t\leq 1\}$ to an element of $\P_k(L)$.
Since isotopies are the identity on the boundary this gives an
injection from the basis of $\P_{k}(L)$ to that of
$\P_{k+1}(L)$.

If there is no source of confusion we will suppress the explicit
dependence on $L$.

\bs{\bf Definition 1.10.} The {\it universal planar algebra}
$\P(L)$ on $L$ is the filtered algebra given by the union of all the
$\P_k$'s ($k=0,1,2,\dots$) with $\P_k$ included in $\P_{k+1}$ as in
the preceding remark.

\bs
A planar algebra will be basically a filtered quotient of $\P(L)$
for some $L$, but in order to reflect the planar structure we need
to impose a condition of {\it annular}
invariance.

\bs{\bf Definition 1.11.} The $j\!-\!k$ {\it annulus} $A_{j,k}$ will
be the complement of the interior of $\b_j$ in
$(j+2)\b_k-(\frac 12,\frac 12)$. So there are $2j$ marked points
on the inner boundary of ${\cal A}_{j,k}$ and $2k$ marked points on
the outer one. An {\it annular} $j-k$ {\it tangle} is the
intersection of a planar network $\cal N$ with $A_{j,k}$ such that
the boundary of $A_{j,k}$ meets $\cal N$ transversally precisely
in the set of marked points of $A_{j,k}$, which are
points on the curves of $\cal N$ other than
endpoints.  The orientation induced by $\cal N$ in
neighborhoods of
$(-\frac 12,-\frac 12)$ and $(0,0)$ are
 required to be positive. Labeling is as
usual.

\bs{\bf Warning.} The diagram in Fig.~1.11(a) is {\it not} an
annular 2--1 tangle, whereas the diagram in Fig.~1.11 (b) is.
\[
        \begin{picture}(0,0)%
\epsfig{file=ps/pic10.ps}%
\end{picture}%
\setlength{\unitlength}{0.00041700in}%
%
\begingroup\makeatletter\ifx\SetFigFont\undefined
% extract first six characters in \fmtname
\def\x#1#2#3#4#5#6#7\relax{\def\x{#1#2#3#4#5#6}}%
\expandafter\x\fmtname xxxxxx\relax \def\y{splain}%
\ifx\x\y   % LaTeX or SliTeX?
\gdef\SetFigFont#1#2#3{%
  \ifnum #1<17\tiny\else \ifnum #1<20\small\else
  \ifnum #1<24\normalsize\else \ifnum #1<29\large\else
  \ifnum #1<34\Large\else \ifnum #1<41\LARGE\else
     \huge\fi\fi\fi\fi\fi\fi
  \csname #3\endcsname}%
\else
\gdef\SetFigFont#1#2#3{\begingroup
  \count@#1\relax \ifnum 25<\count@\count@25\fi
  \def\x{\endgroup\@setsize\SetFigFont{#2pt}}%
  \expandafter\x
    \csname \romannumeral\the\count@ pt\expandafter\endcsname
    \csname @\romannumeral\the\count@ pt\endcsname
  \csname #3\endcsname}%
\fi
\fi\endgroup
\begin{picture}(9368,3330)(871,-5479)
\put(871,-3821){\makebox(0,0)[lb]{\smash{\SetFigFont{12}{14.4}{rm}$(a)$}}}
\put(3001,-3041){\makebox(0,0)[lb]{\smash{\SetFigFont{12}{14.4}{rm}$Q$}}}
\put(8131,-3026){\makebox(0,0)[lb]{\smash{\SetFigFont{12}{14.4}{rm}$Q$}}}
\put(5956,-3806){\makebox(0,0)[lb]{\smash{\SetFigFont{12}{14.4}{rm}$(b)$}}}
\end{picture}

\]
\begin{center}
Figure 1.11
\end{center}

The set of all isotopy classes (isotopies being the identity on the
boundary) of labelled annular $j-k$ tangles, ${\cal A}(L)=\bigcup_{j,k}
{\cal A}_{j,k}(L)$ forms a category whose objects are the sets of
$2j$-marked points of $\b_j$. To compose $A_1\in{\cal A}_{j,k}$ and
$A_2\in{\cal A}_{k,\ell}$, rescale and move $A_1$ so that its outside
boundary coincides with the inside boundary of $A_2$, and the $2k$ 
boundary points match up. Join the strings of $A_1$ to those of $A_2$
at their common boundary and smooth them. Remove that part of the common
boundary that is not strings. Finally rescale the whole annulus so that it
is the standard one. 
The result will depend only on the
isotopy classes of $A_1$ and $A_2$ and defines an element
$A_2A_1$ in ${\cal A}(L)$.

Similarly, an $A\in{\cal A}_{j,k}(L)$ determines a map
$\pi_A:T_j(L)\to T_k(L)$ by surrounding $T\in T_j(L)$ with $A$ and
rescaling. Obviously $\pi_A\pi_B=\pi_{AB}$, and the action of
${\cal A}(L)$ extends to $\P(L)$ by linearity.

\bs
{\bf Definition 1.12} A {\it general planar algebra} will be a
filtered algebra $P=\cup_k P_k$, together with a surjective
homomorphism of filtered algebras, $\Phi:\P(L)\to P$, for some label
set $L$, $\Phi(\P_k)=P_k$, with ker $\Phi$ invariant under ${\cal A}(L)$ 
in the sense that, if $\Phi(x)=0$ for $x\in \P_j$, and
$A\in{\cal A}_{j,k}$ then $\Phi(\pi_A(x))=0$, we say $\Phi$
presents $P$ on $L$.

\bs{\bf Note.} Definition 1.12 ensures that ${\cal A}(L)$ acts on $P$
via $\pi_A(\Phi(x))\eqdef{\rm def}\Phi(\pi_A(x))$.
In particular ${\cal A}(\emptyset)$ ($\emptyset=$ emptyset) acts on
any planar algebra.

\bs
The next results show that this action
extends multilinearly to planar surfaces with several boundary
components.

If $T$ is a planar $k$-tangle (unlabelled), number its boxes
$b_1,b_2,\dots ,b_n$. Then given labelled tangles $T_1,\dots ,T_n$
with $T_i$ having the same number of boundary points as $b_i$, we
may form a labelled planar $k$-tangle $\pi_T(T_1,T_2,\dots ,T_n)$ by
filling each $b_i$ with $T_i$ --- by definition $b_i$ is the image
under a planar diffeomorphism $\th$ of $\b_j$ (for some $j$), and
$T_i$ is in $\b_j$, so replace $b_i$ with $\th(T_i)$ and remove the
boundary (apart from marked points, smoothing the curves at the
marked points). None of this depends on isotopy so the isotopy
class of $T$ defines a multilinear map
$\pi_T: \P_{j_1}\x \P_{j_2}\x\dots\x \P_{j_n}\to \P_k$.
Though easy, the following result is fundamental and its conclusion
is the definition given in the introduction of planar algebras
based on the operad defined by unlabeled planar tangles.

\bigskip
\ni{\bf Proposition 1.13} If $P$ is a general planar algebra
presented on $L$, by $\Phi$, $\pi_T$ defines a multilinear map
$P_{j_1}\x P_{j_2}\x\dots\x P_{j_n}\to P_k$.

{\sc Proof}. It suffices to show that, if all the $T_i$'s but
one, say $i_0$, are fixed in $\P_{j_i}$, then the linear map
$\a:\P_{j_{i_0}}\to \P_k$, induced by $\pi_T$, is zero on ker
$\Phi$. By multilinearity the $T_i$'s can be supposed to be isotopy classes
of labelled tangles. So fill all the boxes other than the $i_0$'th
box with the
$T_i$.  Then we may isotope the resulting picture so that
$b_{i_0}$ is the inside box of a $j_{i_0}-k$ annulus. The map $\a$
is then the map $\pi_A$ for some annular tangle $A$ so
ker $\Phi\subseteq {\op{ker}} \ \pi_A$ by Definition 1.12. \qed

\bigskip
\ni{\bf Proposition 1.14} Let $P$ be a general planar
algebra presented on $L$ by $\Phi$. For each $k$ let $S_k$ be a set
and
$\a:S_k\to P_k$ be a function. Put $S=\coprod_k S_k$. Then there is
a unique filtered algebra homomorphism
$\Theta_S:\P(S)\to P$ with ${\op{ker}} \ \Theta_S$ invariant under
${\cal A}(S)$, intertwining the ${\cal A}(\emptyset)$ actions and such that
$\Theta_S(\tangle{R})=\a(R)$ for $R\in S$.

{\sc Proof}. Let $T$ be a tangle in $P(S)$ with boxes
$b_1,\dots ,b_n$ and let $f(b_i)$ be the label of $b_i$.
We set $\Theta_S(T)=\pi_T(\a(f(b_1)),\a(f(b_2)),\dots ,\a(f(b_n))$
with $\pi_T$ as in 1.13. For the homomorphism property, observe
that $\pi_{T_1T_2}$ and $\pi_{T_1}\dt\pi_{T_2}$ are both
multilinear maps agreeing on a basis. For the annular invariance of
ker $\Theta_S$, note that $\Theta_S$ factors through $\P(L)$, say
$\Theta_S=\Phi\circ \theta$, so that
$\Theta_S(x)=0\Longleftrightarrow
\th(x)\in{\op{ker}} \ \Phi$. Moreover, if $A\in{\cal A}(S)$,
$\theta\!\circ\!\pi_A$ is a linear combination of $\pi_{A'}$'s for
$A'$ in ${\cal A}(L)$. Hence $\Phi(\theta(\pi_A(x))=0$ if
$\theta(x)\in {\op{ker}} \ \Phi$.

Finally we must show that $\Theta_S$ is unique. Suppose we are
given a tangle $T\in \P(S)$. Then we may isotope $T$ so that all its
boxes occur in a vertical stack, as in Figure~1.15.
\[
        \begin{picture}(0,0)%
\epsfig{file=ps/pic11.ps}%
\end{picture}%
\setlength{\unitlength}{0.00041700in}%
%
\begingroup\makeatletter\ifx\SetFigFont\undefined
% extract first six characters in \fmtname
\def\x#1#2#3#4#5#6#7\relax{\def\x{#1#2#3#4#5#6}}%
\expandafter\x\fmtname xxxxxx\relax \def\y{splain}%
\ifx\x\y   % LaTeX or SliTeX?
\gdef\SetFigFont#1#2#3{%
  \ifnum #1<17\tiny\else \ifnum #1<20\small\else
  \ifnum #1<24\normalsize\else \ifnum #1<29\large\else
  \ifnum #1<34\Large\else \ifnum #1<41\LARGE\else
     \huge\fi\fi\fi\fi\fi\fi
  \csname #3\endcsname}%
\else
\gdef\SetFigFont#1#2#3{\begingroup
  \count@#1\relax \ifnum 25<\count@\count@25\fi
  \def\x{\endgroup\@setsize\SetFigFont{#2pt}}%
  \expandafter\x
    \csname \romannumeral\the\count@ pt\expandafter\endcsname
    \csname @\romannumeral\the\count@ pt\endcsname
  \csname #3\endcsname}%
\fi
\fi\endgroup
\begin{picture}(4215,5999)(1794,-6626)
\put(3601,-5761){\makebox(0,0)[lb]{\smash{\SetFigFont{12}{14.4}{rm}$R_1$}}}
\put(3601,-4561){\makebox(0,0)[lb]{\smash{\SetFigFont{12}{14.4}{rm}$R_2$}}}
\put(3526,-2986){\makebox(0,0)[lb]{\smash{\SetFigFont{12}{14.4}{rm}$R_3$}}}
\put(3601,-1786){\makebox(0,0)[lb]{\smash{\SetFigFont{12}{14.4}{rm}$R_4$}}}
\end{picture}

\]
\begin{center}
Figure 1.15
\end{center}

\bs In between each box cut horizontally along a level for which
there are no critical points for the height function along the
curves. Then the tangle becomes a product of single labelled boxes
surrounded by ${\cal A}(\emptyset)$ elements. By introducing kinks as
necessary, as depicted in Figure~1.16, all the surrounded boxes may
be taken in $P_k(S)$ for some large fixed $k$.
\[
        \begin{picture}(0,0)%
\epsfig{file=ps/pic12.ps}%
\end{picture}%
\setlength{\unitlength}{0.00041700in}%
%
\begingroup\makeatletter\ifx\SetFigFont\undefined
% extract first six characters in \fmtname
\def\x#1#2#3#4#5#6#7\relax{\def\x{#1#2#3#4#5#6}}%
\expandafter\x\fmtname xxxxxx\relax \def\y{splain}%
\ifx\x\y   % LaTeX or SliTeX?
\gdef\SetFigFont#1#2#3{%
  \ifnum #1<17\tiny\else \ifnum #1<20\small\else
  \ifnum #1<24\normalsize\else \ifnum #1<29\large\else
  \ifnum #1<34\Large\else \ifnum #1<41\LARGE\else
     \huge\fi\fi\fi\fi\fi\fi
  \csname #3\endcsname}%
\else
\gdef\SetFigFont#1#2#3{\begingroup
  \count@#1\relax \ifnum 25<\count@\count@25\fi
  \def\x{\endgroup\@setsize\SetFigFont{#2pt}}%
  \expandafter\x
    \csname \romannumeral\the\count@ pt\expandafter\endcsname
    \csname @\romannumeral\the\count@ pt\endcsname
  \csname #3\endcsname}%
\fi
\fi\endgroup
\begin{picture}(8948,3521)(1351,-4460)
\put(1351,-1186){\makebox(0,0)[lb]{\smash{\SetFigFont{12}{14.4}{rm}cut}}}
\put(1351,-2461){\makebox(0,0)[lb]{\smash{\SetFigFont{12}{14.4}{rm}cut}}}
\put(6601,-1261){\makebox(0,0)[lb]{\smash{\SetFigFont{12}{14.4}{rm}cut}}}
\put(6601,-2461){\makebox(0,0)[lb]{\smash{\SetFigFont{12}{14.4}{rm}cut}}}
\end{picture}

\]
\begin{center}
Figure 1.16
\end{center}


\bs\ni Since $\Theta_S$ is required to intertwine the $\cal
A(\emptyset)$ action and is an algebra homomorphism, it is determined on
all the surrounded boxes by its value on $\{\tangle{R} : R\in
S\}$, and their products. The beginning and end of $T$ may involve a
change in the value of $k$, but they are represented by an element
of ${\cal A}(\emptyset)$ applied to the product of the surrounded boxes.
So $\Theta_S$ is completely determined on $T$. \qed


\bs{\bf Definition 1.17.} Let $P^1,P^2$  be general planar algebras
presented by $\Phi_1,\Phi_2$ on $L^1,L^2$ respectively.


